
\subsubsection{6.1 Interpretation of Results}

The uniformly zero locality scores and p=1.0 do not constitute evidence that the oracle mechanism is absent. The synthetic data fallback means that no measurement of the oracle mechanism was attempted. Under the available evidence, the oracle/regularizer question remains open.

The DPO training infrastructure ran on synthetic (state, tactic, compiler_error) triples in which all error strings were hardcoded constants. Any probability mass shifts computed during training and evaluation reflect the model's response to meaningless inputs, not responses to real Lean4 compiler semantics. The 100\% state alignment rate reflects alignment of synthetic state IDs, not real proof states.

The gate evaluation correctly recorded a FAIL with p=1.0. This is the only output of the Phase 4 run that accurately reflects the experimental situation.

The hypothesis is not falsified. The Phase 4.5 synthesis document (045_validated_hypothesis.md) classifies the outcome as INFRASTRUCTURE_FAILURE_NOT_SCIENTIFIC and recommends a Phase 4 retry after the LeanDojo environment is correctly installed and the synthetic fallback is removed.

\subsubsection{6.2 The Silent Synthetic Fallback Failure Mode}

The failure mode documented here — a production code path falling back to synthetic data generation when an external tool dependency is absent — poses a risk to any LLM+formal-verifier experiment. The failure is difficult to detect because:

\begin{enumerate}
\item Training runs complete without error.
\item Locality scores are computed and written to result files.
\item Format validation passes.
\item Only the uniformly zero values and subsequent code review reveal the problem.
\end{enumerate}

Plausibility of numerical outputs is not a sufficient validity check. Any pipeline that (a) depends on an external formal tool and (b) has a fallback data generation path in the same codebase faces this risk. This includes pipelines built on Lean4 (LeanDojo), Z3, and static analyzers such as PSL.

The proposed detection protocol is:

```python
real_triples = extract_state_triples(problems[:10], timeout=60)
assert len(real_triples) >= 5, "Pre-run gate FAIL: LeanDojo not producing real triples"
```

This check must be placed before any training begins and must verify that real triples (not synthetic) are returned.

\subsubsection{6.3 Locality Score Metric}

The locality score metric has not yet been validated on real data. Its behavior on real LeanDojo tactic distributions is unknown. The formula includes an ε=1e-8 additive term in the denominator to prevent division by zero. Whether the metric reliably distinguishes oracle from regularizer function on real proof state distributions is an open empirical question.

\subsubsection{6.4 Limitations}

\textbf{L1 — Oracle hypothesis untested.} The central scientific claim — that step-local grounded feedback functions as a local logical oracle — has no direct experimental support from this run. Required fix: install Lean4 toolchain (\texttt{elan} + Lean4 compiler core), verify LeanDojo installation, remove \texttt{_generate_synthetic_triples()} from production code paths, add pre-run gate, and re-execute H-E1.

\textbf{L2 — Vericoding not retrieved.} The Vericoding dataset was not downloaded; 0 problems were loaded. The H-E1 MUST_WORK gate requires valid results only on miniF2F for the PoC stage.

\textbf{L3 — Hard subset not constructed.} Cold-start SFT evaluation on miniF2F (16 rollouts per problem) was not completed, so the hard subset (problems with pass@1 < 20\%) was not identified. DPO training and locality score measurement require the hard subset.

\textbf{L4 — Single seed, 7B scale.} The PoC design uses seed=1 and a 7B parameter model. Statistical validation across seeds and at larger model scale is deferred to H-M3.

\textbf{L5 — Locality score not validated.} The metric has been computed only on synthetic inputs. Its calibration and discriminative power on real proof state distributions are unknown.

\textbf{L6 — Vericoding results are artifacts.} All Vericoding locality scores reported in Table 1 originate from the synthetic data fallback. No actual Vericoding problems were processed.

None of L1–L6 reflects on the scientific validity of the oracle/regularizer hypothesis. All are infrastructure and implementation limitations of this Phase 4 run.

\subsubsection{6.5 Broader Impact}

This research targets LLM theorem proving and formal code verification. The methodological contribution — pre-run environment validation for LLM+formal-verifier pipelines — applies broadly to any AI system that integrates with external formal tools. The silent fallback failure mode identified here could affect safety-critical formal specification tools. Potential negative impacts are limited; this is a methodological paper with no direct deployment pathway.
